\section{国家机构}
\begin{enumerate}
  \item 82 年恢复主席，增设中央军委。18 增加监察委员会。
  \item 体系。人大、人常。一府、两院、一委（监察委）。
  \item 首长负责制。军事和行政机关。其他集体负责。
  \item \zd{全人大}。
        \begin{enumerate}
          \item 上限 3000。特区代表 04 入宪。
          \item 选举权、决定权、法规权（汇总见宪法部分）。
          \item 人事任免。\bs{两院、一委只到县，人大搞正职，人常搞其他}。
                \begin{enumerate}
                  \item \bs{一常、两高、三主席，国家监察委主任}。国家主席、副主席、中央军委主席；最高法、最高检、监察委的老大；人常成员。
                  \item \bs{决定政府和军委}。
                        \begin{enumerate}
                          \item 国家\bs{主席提名} \yjt 国务院\bs{总理，总理提名其他} \yjt 副总理、国务委员（都是国务院二把手） \yjt 部委行署秘正职。
                          \item 军委主席 \yjt 选其他。
                        \end{enumerate}
                  \item \bs{地方检察长，同级人大选，上检提人常}。上级检察长提请同级人常批准。\al{县检察长 \yjt 县人大 \yjt 市检察长 \yjt 市人常}。
                \end{enumerate}
          \item 权利。审批国民经济\bs{计划}和预算、\bs{建省设特；决定战争、和平}。剩余都是人常。
        \end{enumerate}
  \item \zd{全人常}。
        \begin{enumerate}
          \item 4 种人。委员长、副委员长、秘书长（老三，事实上的老二） \yjt 委员长会议。委员（干活的）。
          \item 只有决定 \yjt \zx{只能接收撤职案}。没有选举 \yjt \zx{无权接收罢免}案。即炒鱿鱼方式不一样。
          \item \bs{决定政府和军委，排除国务院总理。两院、一委除正职}。
          \item 权利。
                \begin{enumerate}
                  \item 人常可开火，没说停火。
                  \item 重要类条约。参看三国。\al{全人常让耿爽爽}。
                  \item 特赦。
                  \item 国家勋章、国家荣誉称号、友谊勋章（\al{普京}）。
                \end{enumerate}
        \end{enumerate}
  \item 专委会（智囊团）。\zx{一定是自己人，一定有提案权，一定是主席团提名产生}。
  \item 临时委员会。
        \begin{enumerate}
          \item 调查委员会。
                \begin{enumerate}
                  \item 人员组成。人常 \yjt 能提罢免的人。委员长过五一会议。人大 103 团。
                  \item \zx{可以不公布}调查的情况和材料。
                \end{enumerate}
        \end{enumerate}
  \item \bs{修、修、会、推迟选举}。\zx{修宪、基本法修改、代表开会（所有会）、推迟选举 \yjt 2/3 同意}。必须记住。
  \item （县以上）\zd{人大开会}。
        \begin{enumerate}
          \item \zd{常委会召集，主席团（预备会议选出）主持。乡，老主席团召集，新主席团（开会后现选）主持}。
          \item \zx{召集的人}定时间。
          \item 秘密开会。经过主席团\zx{和}各代表团团长会议决定。
        \end{enumerate}
  \item 议案。3 种。一般（法律议案）、质询、罢免。\bs{机关不质询，人常不罢免。罢免 10、3、团}。
        \begin{enumerate}
          \item 向人大提。
                \begin{enumerate}
                  \item 一般议案。法律法规。
                        \begin{enumerate}
                          \item 机关提案 = \zx{两央、两高、三委（常委、专委、国监委）、一团（主席团）}。
                          \item 个人提案。30 个代表 \hz 1 个代表团。\zx{30、1}。
                        \end{enumerate}
                  \item 质询案。只可能个人提。谴责意味很强。\zx{30、1}。
                  \item \bs{罢免}案。1/10 代表、3 个代表团、主席团。\bs{10、3、团}。
                \end{enumerate}
          \item 人常。\bs{没有一团}。
                \begin{enumerate}
                  \item 机关提案 = 两央、两高、三委（委员长、专委、国监委）。
                  \item 个人提案。1/10 委员长会议、委员。\zd{10}。
                  \item 质询案。\zd{10}。
                  \item \zx{撤职}案。\bs{撤职对象只针对两央，不包括两院、一委}。首长、委员长会议能提。过半。
                \end{enumerate}
          \item \zx{是否列入议程。个人可列可不列，管事的人决定}。人大 \yjt 主席团。人常 \yjt 委员长会议。
          \item 公布。宪法以外，基本都是找\zx{国家主席发主席令}。
        \end{enumerate}
  \item 国家主席。是国家机关。
        \begin{enumerate}
          \item 年龄。特首 58 同乘，主席 95 至尊。
          \item \bs{吃饭}（国事活动）、\bs{请吃饭、发友谊勋章}（两央、一委提，常委会定）。
        \end{enumerate}
  \item 国务院。性质中央政府。最高行政机关。
        \begin{enumerate}
          \item 总理 \yjt 副总理、国务委员 \yjt 三级秘书长。
          \item 组成。\zx{包括发改委，不包括国资委}（直属特设机构）。
          \item 老大对人大人常负责、领导工作、提名成员、最后决定、最后签字（机关令）。
          \item \zd{部委调整，总理提，人大人常定}。\bs{一高一低留中间}。
          \item 双领导则是双负责。如\zx{审计、公安、检察}。
          \item \zd{不正常}状态。战争、动员、宣战、紧急状态。\bs{人大（开停火）、人常}、\bs{国务院（省内部分地区、紧急）}。
        \end{enumerate}
  \item 中央军委。\zx{不报告工作，也不接受质询}。
  \item 有连任限制的。\bs{人常、政府正副职，两院、一委正加特}。正职加特首。
  \item 监察委员会。组成、负责制、领导体制。
  \item 地方国家机关。5 级人大，4 级常委专委，3 级秘书长。
        \begin{enumerate}
          \item 人事任命。
          \item 特殊任免。直辖市、地区\zx{中院和分检，人常}（没有对应人大，提高半级）。
                \begin{enumerate}
                  \item \zx{法院。人常会议提名，人常决定 \yjt 院长 \yjt 副院长等}。
                  \item 检察院。\zx{检察长整体打包提名，人常决定}。
                \end{enumerate}
          \item 会议制度。
                \begin{enumerate}
                  \item 人常 2 月开一次会。主任会议。
                  \item 乡人大每年两回。出席人数统一 2/3。
                \end{enumerate}
          \item 议案。
          \item 撤职。\zd{人常 \yjt 一府、两院除正职，特殊法检含正职}。\zx{谁任命谁来撤}。
          \item 部委调整。上一级政府定，报同级人常。\bs{本级提，上级批，同人备}。
          \item 协同。立法（地级市以上）、行政协同（县级以上）= 部门协同（共同上级）+ 区域协同。
        \end{enumerate}
\end{enumerate}